\subsection{格序群}

回忆，其中每个非空有限子集都有上确界和下确界的有序集称为格（E，III，第13页）。显然，格序群的乘积，特别地全序群的乘积，是格序群。与此相反，格序群的子群不一定是格序的。

于是，在乘积序群 \(\mathbf{Z} \times \mathbf{Z}\) 中，“反对角线”（数对 \((n, n')\) ——满足 \(n + n' = 0\) ——组成的集合）装备离散序，因此不是格序群。* 单个实变量的多项式加法群（VI，第1页，例2）是滤过群（因为 \(p(x)\) 与 \(q(x)\) 都小于 \((p(x))^2 + (q(x))^2 + 1\) ），可以证明它不是格序的。*

命题7. --- 若 \(x\) 与 \(y\) 是序群 \(G\) 的两个元素，且元素 \(\inf(x, y)\) , \(\sup(x, y)\) 之一存在，则另一个也存在，并且 \(x + y = \inf(x, y) + \sup(x, y)\) 。

事实上，根据关系式(1)和(2)（VI，第9页），我们有

\[
\sup (a - x, a - y) = a + \sup (- x, - y) = a - \inf (x, y),
\]

只需取 \(a = x + y\) 即可。

命题7（DIV）。 --- 若 \(a, b \in \mathbf{K}^{*}\) ，且若 \(d\) 是 \(a\) 与 \(b\) 的一个最大公因子，而 \(m\) 是 \(a\) 与 \(b\) 的一个最小公倍元，则乘积 dm 是 \(ab\) 的相伴元。

命题 8. --- 设P是序群G的正元素集合。要使G是格序群，必要且充分的是G = P - P，且此外P（装备诱导序）满足下列两个条件之一：

\begin{enumerate}
\def\labelenumi{\alph{enumi})}
\item
  \(\mathbf{P}\) 的每一对元素在 \(\mathbf{P}\) 中都有上确界。
\item
  \(\mathbf{P}\) 的每一对元素在 \(\mathbf{P}\) 中都有下确界。
\end{enumerate}

这些条件的必要性是显然的：事实上，关系 \(G = P - P\) 表明G是滤过的（VI，第4页，命题4）；另一方面，P的两个元素在G中的上确界和下确界是正的，因而它们也是这两个元素在P中的上确界和下确界。

反过来，首先注意到在假设a)（相应地，b)）下，P中任意一对元素x、y在G中的上确界（相应地，下确界）等于其在P中的上确界a（相应地，下确界b）。这对a)是显然的，因为x和y的任何上界都是正的；对于b)，设 \(z \in G\) 是 x 和 y 的一个下界；则存在 \(u \in P\) 使得 \(z + u \in P\) ，因为 G = P - P；现在 \(\inf_{P}(x + u, y + u)\) 大于或等于 \(b + u\) ，因此具有形式 \(b + c + u\) ( \(c \geqslant 0\) ）；因为 \(b + c\) 小于x且小于y，我们有 c = 0；因此 \(\inf_{P}(x + u, y + u) = b + u\) ，这意味着 \(z + u \leqslant b + u\) ，因此 \(z \leqslant b\) 且 b 显然是 x 和 y 在 G 中的下确界。现在若 x 和 y 是 G 的任意元素，我们将它们平移进 P：令 \(v \in P\) 使得 \(x + v\) 与 \(y + v\) 都是正的（VI，第 5 页，命题 5）；在假设 a)（相应地 b)）下

\(x + v\) 与 \(y + v\) 在 P 中存在上确界（相应地，下确界），因而根据刚才所证，在 G 中也存在；通过平移，x 和 y 在 G 中有上确界（相应地，下确界）；由命题7，其中一个的存在蕴含另一个的存在，这就表明这些条件是充分的。

